\documentclass[11pt]{article}

\usepackage[a4paper,margin=1in]{geometry}
\usepackage{booktabs}
\usepackage{tabularx}
\usepackage{amsmath}
\usepackage{amssymb}
\usepackage{enumitem}
\usepackage{xcolor}
\usepackage[hidelinks]{hyperref}

\newcommand{\pdeta}{\textsf{pDetA}}
\newcommand{\passa}{\textsf{pAssA}}
\newcommand{\phota}{\textsf{pHOTA}}
\newcolumntype{Y}{>{\raggedright\arraybackslash}X}

\title{\textbf{Predicting Zero-Shot Video-Tracker Transfer:\\ Methods, Results, and Engineering Log}}
\author{Progress report --- ``will-it-track''}
\date{\today}

\begin{document}
\maketitle

\begin{abstract}
\noindent We are testing whether the zero-shot tracking performance of a frozen promptable
video model (SAM~3) on an \emph{unseen} species or place can be predicted \emph{before running it},
from four label-free ``distances'' (taxonomic, visual, environment, temporal), and whether the two
halves of tracking --- \emph{finding} the animal (\pdeta{}) and \emph{following} it (\passa{}) --- are
governed by different factors. This document records, in plain language, exactly what has been built and
run so far, what the results say, and every practical and statistical problem we hit together with the
fix we applied. The headline so far is an \emph{honest null}: on the two splits tested, the distances do
not predict SAM~3's per-cell detection/association to a degree any valid test supports. This is a
legitimate scientific result, and it points clearly at the next experiment (a powered species hold-out),
which is running now.
\end{abstract}

\section{Aim and hypotheses}
Conservation teams point promptable trackers at new species and new camera-trap sites constantly, but
nobody can say in advance whether the model is trustworthy for a given animal in a given place. We try to
fill that gap with a \emph{reliability estimator}: given only label-free distances, predict \pdeta{} and
\passa{} with honest confidence intervals. The guiding hypotheses are:
\begin{description}[leftmargin=2.2em,itemsep=2pt]
  \item[H1 (detection).] \pdeta{} falls with \textbf{species novelty} (taxonomic + visual distance).
  \item[H2 (association).] \passa{} falls with \textbf{environment difficulty} (scene distance, day/night/IR,
        clutter), largely regardless of species.
  \item[H0 (null).] Distance explains little variance --- \emph{still a valid result}, and one that would
        redirect the project toward representational probing.
\end{description}
The goal is explicitly \textbf{not} to make the tracker score higher; it is to \emph{predict} and
\emph{explain} its transfer.

\section{Data}
We use \textbf{SA-FARI} (Meta $\times$ Conservation X Labs, 2025), the largest open multi-animal tracking
dataset: 11{,}609 camera-trap videos, 99 species with a full 7-level taxonomy, 741 locations on 4
continents, and per-video timestamps. Annotations are on 6\,fps frames (mean $\approx$85 frames per clip).
Crucially, SAM~3 is \textbf{frozen and zero-shot} on all of SA-FARI, so the train/test split is only our
own reference anchor for measuring distances --- not anything the model learned.

\section{Method}
The pipeline has five stages; each is a small, config-driven module.

\subsection{Inference}
A frozen SAM~3 promptable tracker is run over every probe of a split. For each video we give a text prompt
(the species name, e.g.\ \texttt{"impala"}) and the model segments and tracks every matching animal across
the clip. Predictions are written one JSON per video, so the run is \emph{resumable}. Hard negatives
(a query for an absent species) are kept --- they should return nothing.

\subsection{Scoring}
We score with the \emph{official} SA-Co / VEval evaluator (never re-implemented). The promptable HOTA
score decomposes cleanly into the two parts we analyse separately:
\[
  \phota{} \;=\; \sqrt{\pdeta{} \times \passa{}}, \qquad
  \pdeta{} = \text{did it find the animals?}, \quad
  \passa{} = \text{did it keep each identity over time?}
\]
Scores are recorded per \textbf{cell} --- a (species, location, time) group --- together with the
\emph{support} behind each score (number of frames, masklets, videos).

\subsection{The four label-free distances}
Every distance is computable \emph{without any label of the target species}, which is what makes a
before-running prediction possible.

\begin{table}[h]
\centering
\small
\begin{tabularx}{\textwidth}{lYl}
\toprule
Distance & What it measures & Granularity \\
\midrule
Taxonomic   & Tree distance (Kingdom$\to$Species) to the nearest reference species & per species \\
Visual      & $1-\cos$ between DINOv2 embeddings of the \emph{mask-cropped animal} and the nearest
              reference species' prototype & per species \\
Environment & $1-\cos$ between DINOv2 embeddings of the \emph{masked-background scene} and the nearest
              reference location's prototype; plus a night/IR flag and a clutter score & per location \\
Temporal    & Gap (in years) between the probe footage and the reference footage & per cell \\
\bottomrule
\end{tabularx}
\caption{The four distances. Visual and environment reuse a shared, cache-first embedding pipeline.}
\end{table}

\subsection{The two experiments (splits)}
\begin{description}[leftmargin=1.4em,itemsep=3pt]
  \item[Split B --- location hold-out (secondary, tests H2).] Reference $=$ official train split,
        probe $=$ test split. Camera sites are unseen; species are largely shared, so the taxonomic
        distance is $\approx 0$ and the environment/temporal distances vary. This is the ``does the scene
        matter?'' experiment.
  \item[Split A --- species hold-out (primary, tests H1).] Leave-\emph{one}-species-out over the pooled
        present set: each probe species' distance is to the nearest \emph{other} species, which makes
        taxonomic novelty a genuine \emph{graded} axis. This is the ``does novelty matter?'' experiment.
\end{description}

\subsection{Modelling}
For each target (\pdeta{}, \passa{}) we fit a support-weighted logit-link GLM (a fractional-logit model
on the bounded score), with the four distances plus the scene covariates and a $\log(\text{support})$
covariate so that ``rare'' is never mistaken for ``far''. Predictors are standardised. We then:
\begin{itemize}[leftmargin=1.4em,itemsep=1pt]
  \item report \textbf{group-cluster-bootstrap} confidence intervals (refit while resampling whole
        species / locations) --- see the confidence-interval issue below;
  \item run \textbf{grouped cross-validation} (leave-one-species-out and leave-one-location-out) and test
        the out-of-sample gain over a mean predictor with a \textbf{paired group-bootstrap} $p$-value;
  \item partition variance across factors (dominance analysis) and check collinearity (VIF).
\end{itemize}

\section{Results so far}

\subsection{Gate 1 --- measurement (is SAM~3 behaving sensibly?)}
Zero-shot SAM~3 on the SA-FARI test split reproduces the published ballpark. The dataset-level (micro)
mask \phota{} is $\approx 0.65$ with a real detection--association gap. Table~\ref{tab:gate1} reports the
per-cell (macro) support-weighted means we feed to the models; these are lower than the micro number
because 78\% of cells are single-object clips where \passa{}$\approx$\pdeta{} (finding and following
coincide when there is only one animal).

\begin{table}[h]
\centering
\begin{tabular}{lr}
\toprule
Metric (mask, per-cell macro mean) & Value \\
\midrule
\pdeta{} & 0.527 \\
\passa{} & 0.546 \\
\phota{} & 0.533 \\
\midrule
Cells scored & 346 of 1447 \\
\bottomrule
\end{tabular}
\caption{Gate 1 measurement on the test split. The 1101 unscored cells are hard negatives (species
absent, no ground-truth track), correctly excluded from the HOTA fit and reserved for a separate
false-positive analysis.}
\label{tab:gate1}
\end{table}

\subsection{Gate 2 --- location split (Split B): a null}
After correcting the inference (the confidence-interval issue below), \textbf{every distance coefficient's confidence
interval spans zero}, and the out-of-sample gain over a mean predictor is \textbf{not significant} on
either cross-validation scheme (Tables~\ref{tab:coefB}--\ref{tab:cv}). The large-looking taxonomic
coefficient ($-0.36$) is a fragile artifact of this split (see the novelty-axis issue below), not a graded
effect.

\begin{table}[h]
\centering
\small
\begin{tabular}{lcc}
\toprule
Factor & \pdeta{} coef.\ [95\% CI] & \passa{} coef.\ [95\% CI] \\
\midrule
Taxonomic distance   & $-0.362$ $[-3.87,\,0.09]$ & $-0.368$ $[-3.40,\,0.08]$ \\
Visual distance      & $-0.144$ $[-0.50,\,0.14]$ & $-0.122$ $[-0.49,\,0.18]$ \\
Environment distance & $+0.042$ $[-0.19,\,0.24]$ & $+0.013$ $[-0.23,\,0.21]$ \\
Clutter              & $-0.066$ $[-0.38,\,0.29]$ & $-0.087$ $[-0.41,\,0.28]$ \\
Night / IR           & $-0.405$ $[-0.94,\,0.11]$ & $-0.462$ $[-1.02,\,0.11]$ \\
$\log(\text{support})$ & $-0.216$ $[-0.51,\,0.17]$ & $-0.224$ $[-0.53,\,0.18]$ \\
\bottomrule
\end{tabular}
\caption{Split B (location) standardised coefficients with group-cluster-bootstrap CIs. Every distance
spans zero. The temporal gap is omitted: it is constant on this split and carries no information.}
\label{tab:coefB}
\end{table}

\subsection{Gate 2 --- species hold-out (Split A, preview): a cleaner null}
Because the location split cannot properly test novelty, we ran the species hold-out (leave-one-species-out)
on the already-scored test species --- a free preview that needs no new inference. With a \emph{graded}
novelty axis, the taxonomic effect \textbf{collapses to essentially zero} (Table~\ref{tab:coefA}),
confirming that the location split's $-0.36$ was an artifact. The out-of-sample gain is again not
significant --- indeed slightly \emph{negative} (the distance model is no better than, or marginally worse
than, predicting the mean).

\begin{table}[h]
\centering
\small
\begin{tabular}{lcc}
\toprule
Factor & \pdeta{} coef.\ [95\% CI] & \passa{} coef.\ [95\% CI] \\
\midrule
Taxonomic distance   & $+0.026$ $[-0.31,\,0.31]$ & $+0.008$ $[-0.33,\,0.31]$ \\
Visual distance      & $+0.120$ $[-0.23,\,0.38]$ & $+0.135$ $[-0.23,\,0.40]$ \\
Environment distance & $-0.025$ $[-0.30,\,0.21]$ & $-0.046$ $[-0.33,\,0.19]$ \\
Clutter              & $-0.106$ $[-0.37,\,0.21]$ & $-0.135$ $[-0.41,\,0.20]$ \\
Night / IR           & $-0.394$ $[-0.99,\,0.09]$ & $-0.466$ $[-1.09,\,0.06]$ \\
$\log(\text{support})$ & $-0.188$ $[-0.49,\,0.19]$ & $-0.193$ $[-0.51,\,0.21]$ \\
\bottomrule
\end{tabular}
\caption{Split A (species hold-out, test-species pool) standardised coefficients. With a graded novelty
axis the taxonomic and visual effects are $\approx 0$ and span zero. pseudo-$R^2$ is 0.033 / 0.039.}
\label{tab:coefA}
\end{table}

\begin{table}[h]
\centering
\small
\begin{tabular}{lllrrr}
\toprule
Split & Hold-out & Target & $\Delta$MAE & 95\% CI & $p$ \\
\midrule
Location (B) & location & \passa{} & $+0.007$ & $[-0.005,\,0.018]$ & 0.13 \\
Location (B) & location & \pdeta{} & $+0.004$ & $[-0.007,\,0.015]$ & 0.23 \\
Species (A)  & species  & \passa{} & $-0.003$ & $[-0.013,\,0.006]$ & 0.74 \\
Species (A)  & species  & \pdeta{} & $-0.004$ & $[-0.013,\,0.005]$ & 0.79 \\
\bottomrule
\end{tabular}
\caption{Out-of-sample gain over a mean predictor ($\Delta$MAE $>0$ means the model is better), with a
paired group-bootstrap CI and one-sided $p$. \textbf{No row is significant} on either split.}
\label{tab:cv}
\end{table}

\subsection{What the results mean (honestly)}
On both splits tested, the label-free distances do \textbf{not} predict SAM~3's per-cell detection or
association to a degree any valid, group-clustered test supports. This is \textbf{H0-leaning}. It is not
yet ``H0 confirmed for the whole thesis'': the species hold-out preview uses only the 83 test species, and
the powered version (adding a sample of train species) is running now. The two \emph{large-looking}
effects that survive as point estimates --- night/IR and (on Split B) taxonomic --- do not survive an
honest interval, and the taxonomic one is a known artifact. The one robust qualitative signal is that
\textbf{night/IR footage lowers both} \pdeta{} and \passa{} (day $\to$ night: \pdeta{} $0.60\to0.51$,
\passa{} $0.62\to0.52$), but even this is not individually significant once clustering is respected.

\section{The evidence for the null (H0)}
H0 is the hypothesis that \emph{distance explains little variance} --- that the four label-free distances
carry almost no usable information about how well SAM~3 will track a new species in a new place. This is
not a fallback we reached for; it is what six independent lines of evidence point to. Below, each line is
stated, then explained in plain terms, then linked back to H0.

\paragraph{How to read the numbers (a short primer).}
A \emph{coefficient} is the slope of a fitted line: how much the score is predicted to change when a
distance goes up by one standard deviation. A \emph{95\% confidence interval (CI)} is the range of slopes
the data are consistent with; if that range \emph{includes zero}, we cannot rule out that the true effect
is \emph{nothing}. A \emph{$p$-value} is the probability of seeing an effect at least this large if the
truth were really zero --- small $p$ (below $0.05$) means ``unlikely to be a fluke''; large $p$ means
``easily just noise''. \emph{Pseudo-$R^2$} is a 0--1 measure of how much of the variation the model
explains (0 $=$ nothing, 1 $=$ everything). \emph{MAE} (mean absolute error) is the average size of the
prediction mistakes; the \emph{baseline} is the MAE you get by ignoring the distances and just guessing the
overall average score. \emph{Out-of-sample} means the model is judged on species/locations it never saw
during fitting --- the only fair test of a predictor.

\subsection{The six lines of evidence}
\begin{enumerate}[leftmargin=1.6em,itemsep=6pt]
  \item \textbf{The models explain almost none of the variation.}
        The fitted models reach a pseudo-$R^2$ of only $\approx 0.055$ on the location split and
        $\approx 0.03$--$0.04$ on the species hold-out (Table~\ref{tab:coefA} caption). In words: even
        \emph{inside} the data it was fitted on, the model captures around 3--6\% of what makes one cell's
        score differ from another's --- and most of even that little is carried by the night/IR flag and
        the support covariate, not by the novelty/environment distances the thesis is about. A predictor
        that explains 3--6\% of the variance is, for practical purposes, explaining nothing.

  \item \textbf{Every distance's confidence interval includes zero.}
        After honest inference (resampling whole species/locations --- see the confidence-interval issue
        in the log), not one of the four distances has a coefficient we can distinguish from zero, on
        either split or either target (Tables~\ref{tab:coefB} and~\ref{tab:coefA}). For example the
        taxonomic slope on the species hold-out is $+0.026$ with a CI of $[-0.31, +0.31]$: the data are
        equally consistent with a small positive effect, a small negative effect, and no effect at all.
        When every interval straddles zero, the data are telling us they cannot pin down any direction ---
        the signature of H0.

  \item \textbf{Even the raw, un-modelled correlations are tiny.}
        Before any modelling, the plain correlation between each distance and the score never exceeds
        $|r| \le 0.17$, i.e.\ $r^2 \le 0.03$ (at most 3\% of variance shared). This matters because it
        rules out the objection ``maybe the model is just badly specified'': the weakness is already there
        in the simplest possible look at the data, not introduced by the GLM.

  \item \textbf{The predictor fails its own out-of-sample test --- the decisive one.}
        The whole point was a \emph{predictor}: distances in, a useful guess of the score out, for
        species and places never seen. We tested exactly that with grouped cross-validation (hold out a
        whole species or a whole location, fit on the rest, predict the held-out group). The distance
        model does \textbf{not} significantly beat the trivial ``just guess the average'' baseline
        (Table~\ref{tab:cv}): on the location split the gain is a mere $0.004$--$0.007$ MAE with
        $p = 0.13$--$0.23$ (not significant), and on the species hold-out the gain is actually
        \emph{negative} ($p = 0.74$--$0.79$) --- the distances make the out-of-sample guess slightly
        \emph{worse} than ignoring them. A predictor that cannot beat ``guess the mean'' has no practical
        value, and this is the single strongest piece of evidence because it is the deliverable failing
        its actual job.

  \item \textbf{The one strong-looking effect was an artifact that vanished under a cleaner design.}
        On the location split the taxonomic distance looked important ($-0.36$). But that split makes
        taxonomic novelty degenerate: almost all test species also occur in the reference set, so their
        distance is $\approx 0$, and only $\sim$7 genuinely-unseen species sit above zero --- the ``$-0.36$''
        was a fragile contrast driven by those few. When we re-ran with a proper \emph{graded} novelty
        axis (the species hold-out, where each species is compared to the nearest \emph{other} species),
        the taxonomic effect collapsed to $+0.03$/$+0.01$ (Table~\ref{tab:coefA}). An apparent signal that
        evaporates the moment the design is fixed is exactly what you expect when there is no real effect.

  \item \textbf{The null is consistent across two independent designs and both targets.}
        We now have four independent looks --- \{location split, species split\} $\times$ \{detection,
        association\} --- and all four give the same answer: no significant distance effect, no
        out-of-sample gain. A genuine effect would tend to surface in at least one of them; four
        consistent nulls, from two differently-constructed splits, make a coincidental miss unlikely and
        strengthen the H0 reading.
\end{enumerate}

\subsection{What this does, and does not, establish}
H0 here is a \emph{scoped} claim, and it is important to say exactly how far it reaches.
\begin{itemize}[leftmargin=1.4em,itemsep=2pt]
  \item \textbf{It does establish:} these four label-free distances, at the \emph{per-cell} granularity,
        on the two splits measured, do not predict SAM~3's detection or association to any degree a valid
        test supports.
  \item \textbf{It does not establish} that \emph{no} label-free signal exists anywhere. Three things are
        untested and could still carry signal: (i) coarser aggregation (per-species or per-location
        averages, which are far less noisy than single short clips); (ii) a \emph{familiarity proxy}
        (separability in SAM~3's own feature space, rather than in DINOv2's); and (iii) the
        false-positive / hallucination regime --- whether a novel-species prompt fires on an
        \emph{absent} species --- which our present analysis, conditioned on the animal being present,
        never touches.
  \item \textbf{A residual qualitative pattern remains:} night/IR footage lowers both scores
        (day $\to$ night: \pdeta{} $0.60 \to 0.51$, \passa{} $0.62 \to 0.52$). Environment therefore does
        matter a little --- but through a simple day/night flag, not through the continuous scene-distance
        the hypothesis proposed, and even this is not individually significant once clustering is respected.
\end{itemize}

\subsection{Why a null is still a result --- and where it points}
A validated negative result on a genuinely new problem (the first before-running transfer predictor for a
promptable video \emph{tracker}, split into detection vs association) is itself a contribution: it tells
the field that the ``obvious'' label-free distances are not enough, and it does so with honest,
cluster-aware statistics rather than an over-confident claim. It also \emph{redirects} the project. If
transfer is not simply a function of how \emph{far} a new species/scene is, the sharper question becomes
whether SAM~3 \emph{internally represents} the structure at all --- \textbf{representational probing}:
instead of ``does distance predict the score?'', ask ``do SAM~3's own features linearly encode taxonomy,
novelty, or scene type?''. That is a mechanism question, it reuses the same frozen model and data, and it
is exactly the pivot H0 was written to motivate.

\section{Issues encountered and how we solved them}
This section is the engineering and statistics log. Each row is a real problem we hit and the fix applied.

\subsection{Engineering / infrastructure}
\begin{table}[h]
\centering
\small
\begin{tabularx}{\textwidth}{p{4.3cm}Y}
\toprule
Issue & Solution applied \\
\midrule
Feature assembly ran out of memory at full sampling caps (all cropped frames held in RAM at once on a
$\sim$58\,GB container). &
Rewrote the embedding step to work in \emph{chunks}: load $\to$ embed $\to$ cache $\to$ free, so peak RAM
is bounded by one chunk regardless of the caps. RSS stayed at $\sim$7\,GB. \\
\addlinespace
Embedding was I/O-bound at $\sim$5 crops/s (the GPU sat idle while single-threaded frame reads from a
network filesystem blocked). &
Parallelised frame load/crop with a thread pool ($\sim$32 workers); the GPU is now the bottleneck. \\
\addlinespace
The official scorer failed on a fresh pod (its import-time dependencies were missing). &
Hardened the pod setup to clone the scorer and install its undeclared deps
(\texttt{iopath, einops, timm, ftfy, regex}). \\
\addlinespace
PyTorch pulled a CUDA build newer than the pod's driver; and \texttt{pip} was OOM-killed installing torch
(the temp dir was RAM-backed). &
Pin torch to the driver's CUDA; move \texttt{TMPDIR} onto persistent disk and install with
\texttt{-{}-no-cache-dir}. \\
\addlinespace
The full train-split inference is $\sim$31{,}000 probes $\approx$ \textbf{10 days} of GPU time --- far too
expensive for one experiment. &
Added a \emph{per-species cap} (a stratified sample of $\sim$30 positive videos per species, $\approx$1500
probes $\approx$6\,h) that still gives a graded novelty axis. \\
\addlinespace
A single pathological clip crashed SAM~3's forward pass (a likely CUDA out-of-memory on a long video) and
killed the entire multi-hour batch. &
Wrapped per-video inference in a guard: a crashing clip is \emph{skipped} (scored as a miss), the GPU
cache is freed, and the batch continues. \\
\addlinespace
Long clips could exhaust GPU memory during video preprocessing. &
Keep SAM~3 video preprocessing on the CPU (raw storage only), running the model itself on the GPU. \\
\addlinespace
The proxy SSH to the pod allows no \texttt{scp} and needs a pseudo-terminal. &
Transfer files by \texttt{tar}$+$\texttt{base64} over the terminal and decode locally; run the analysis
parquets locally where possible. \\
\bottomrule
\end{tabularx}
\end{table}

\subsection{Statistical / scientific}
\begin{table}[h]
\centering
\small
\begin{tabularx}{\textwidth}{p{4.3cm}Y}
\toprule
Issue & Solution applied (or proposed) \\
\midrule
\textbf{Over-precise confidence intervals.} Weighting each cell by its frame count
inflated the effective sample $\sim$85$\times$ and shrank every standard error $\sim$9$\times$; on top of
that, predictors are constant within a species/location (pseudo-replication), so 346 cells are really only
$\sim$60--92 independent groups. The naive intervals looked ``significant'' when they were not. &
Replaced the naive intervals with a \textbf{group-cluster-bootstrap}: resample whole species / locations,
refit, take percentile intervals. Every distance CI then spans zero. \emph{Next:} also report
cluster-robust (sandwich) SEs as a cross-check. \\
\addlinespace
\textbf{No significance test on the cross-validation gain.} A $+0.005$ MAE edge over the baseline was being
read as ``beats baseline'' with no test. &
Added a \textbf{paired group-bootstrap} on the CV gain (resample whole held-out groups) $\to$ a CI and
one-sided $p$. All gains are non-significant. \\
\addlinespace
\textbf{Degenerate temporal feature.} On the location split the temporal gap is present for only 32/346
cells and constant among them --- an all-zero column after standardising. &
Drop zero-variance predictors automatically. Temporal distance is therefore \emph{untested} (not falsified)
on this split; it needs a time-based split to test. \\
\addlinespace
\textbf{Novelty axis was degenerate on Split B.} With the train split as reference,
the taxonomic distance is $\approx 0$ for shared species and $>0$ for only $\sim$7 truly-unseen ones, so
the ``$-0.36$'' coefficient was a fragile 7-species spike-at-zero, not a graded effect. &
Run \textbf{Split A} (leave-one-species-out), where novelty is graded. The taxonomic effect then collapses
to $\approx 0$ --- confirming the artifact. \\
\addlinespace
\textbf{Environment sign-flip.} The continuous scene distance has a weak \emph{negative} raw correlation
but flips to $\approx 0$/positive in the fit --- a suppression symptom of its collinearity with the
night/IR flag (both are per-location, from the same background). &
Flagged; VIF currently covers only the four distances. \emph{Next:} add night/IR and clutter to the VIF
check and refit environment with/without night/IR to show the confound explicitly. \\
\addlinespace
\textbf{Detection-vs-association is under-identified.} 78\% of cells are single-object, so \passa{}$\approx$
\pdeta{} and the two regressions are near-duplicates --- the project's core decomposition cannot be tested
here. &
\emph{Next:} a targeted analysis on the $\sim$22\% multi-object cells, where association carries information
independent of detection (flagged exploratory, low power). \\
\addlinespace
\textbf{Hard negatives excluded from the fit.} The detection model is fit on animal-present cells only
(recall-given-present), so the false-positive / hallucination regime is not yet covered. &
\emph{Next:} a separate false-positive-rate-vs-novelty analysis that reintroduces the 1101 hard negatives. \\
\addlinespace
\textbf{Scope of the novelty claim.} ``Distance to the SA-FARI train reference'' is not distance to SAM~3's
true (unknown) pretraining corpus. &
Scope every claim to the train split, and report a \emph{familiarity proxy} (separability in SAM~3's own
feature space) alongside. \emph{Next:} implement the proxy (currently deferred). \\
\bottomrule
\end{tabularx}
\end{table}

\section{Current status and next steps}
\begin{itemize}[leftmargin=1.4em,itemsep=1pt]
  \item \textbf{Done and validated:} the full measurement + modelling pipeline; the location split (Split B,
        null); the species hold-out preview (Split A on test species, cleaner null); honest clustered
        inference throughout.
  \item \textbf{Running now:} the \emph{powered} species hold-out --- a capped ($\sim$30 videos/species)
        SAM~3 pass over the train split ($\sim$1500 probes, $\sim$6\,h), which will add species and cells to
        tighten the H1 test.
  \item \textbf{Then, in priority order:} (i) refit Split A with the pooled train+test species; (ii) a
        multi-object-only association analysis; (iii) the hard-negative false-positive analysis;
        (iv) the familiarity proxy; (v) if the powered test confirms the null, pivot the thesis toward
        \emph{representational probing} (does SAM~3 encode taxonomy?), which H0 explicitly motivates.
\end{itemize}

\paragraph{Bottom line.} The engineering works end-to-end and the statistics are now honest. The current,
carefully-tested answer is that simple label-free distances do not predict this tracker's zero-shot
transfer on the splits measured so far --- a clean, publishable negative result that sharpens the next
question rather than a failure.

\end{document}
